\documentclass[12pt]{article}
\usepackage[margin=1in]{geometry}
\setlength{\parskip}{5pt}
\usepackage{amsmath}
\usepackage{mathtools}
\usepackage{graphicx}


\title{ECE 496 - Silicon Integrated Circuit Design Technology  Lab 1 - Cleanroom Safety and Tour}
\author{Gianni Avilan-Losee}
\date{\today}
\begin{document}
\maketitle
\section{Objectives}

This lab's objective was to tour the cleanroom and familiarize ourselves with safety precautions, gowning procedure, the layout of the lab, and the equipment we will be using.

\section{Background}
The first important safety step is to dress appropriately for the environment. In our case, it's recommended that we use long pants, close-toed shoes, and long-sleeved shirts. Essentially, we want to aim for as little exposed skin as possible \textit{before} gowning. 

Situational awareness is vital, ensuring that we do not knock over or trip over any equipment or dangerous substances. We should know at all times where the emergency exits are, as well as being generally familiar with the cleanroom's layout, along with fire extinguisher locations and first aid kit locations. Other important equipment to be aware of is phones, safety shower/eyewash, chemical spill kit, and circuit breaker panel.

\section{Procedure}
In the cleanroom, we began by following the gowning procedure, which involved applying shoe covers, bouffant, hood, face mask, booties, cleanroom gloves, and safety eyewear.

In lab, we were given a tour of the three bays in the cleanroom: the etch and oxide deposition lab, the photolithography lab, and the deposition and characterization lab. 

In the first bay, the equipment included an oxidation and diffusion furnace, used to grow a layer of silicon dioxide onto the silicon wafer and implant impurities, such as boron or arsenic, into the silicon. This bay also includes a hood, as well as a spin rinse dryer and wet bench.  

In the second bay, the equipment included a spin coater and mask aligner, used for depositing even layers of liquid material onto wafers and UV photolithographic etching, respectively. The bay also includes a measuring microscope, hot plates, and a hood. 

In the third bay, the equipment included a Nano36 Sputter System and CHA e-beam evaporator, both used for depositing dielectric films or metals onto wafers. 

\section{Results}

Not applicable for this lab.

\section{Conclusion}

In conclusion, we have learned proper safety procedures to follow in lab, as well as appropriate use of PPE, and the function of various equipment found in the lab, along with the layout and locations of first aid and other emergency equipment.

\end{document}

